
% https://tex.stackexchange.com/a/59706/295527

\documentclass{article}
%\usepackage[T1]{fontenc}
\usepackage{ntheorem}
\newtheorem{theorem}{Theorem}
\usepackage{amsmath}
\DeclareMathOperator{\Res}{Res}

%% Choose one of the following (if not choosing the  
%% default, viz., Computer Modern, font family):
%\usepackage{lmodern}
%%
%\usepackage{mathpazo}
%\usepackage[theoremfont]{newpxmath} \usepackage{newpxmath}
%\usepackage{kpfonts}
%%
%\usepackage{mathptmx}
%\usepackage{times,mtpro2}
%\usepackage{stix}
%\usepackage{txfonts}
%\usepackage{newtxtext,newtxmath}
%%
%\usepackage{libertine} \usepackage[libertine]{newtxmath}
%\usepackage{libertine,libertinust1math} % added 2019/11/28
%%
%\usepackage{newpxtext} \usepackage[euler-digits]{eulervm}

\begin{document}\pagestyle{empty}
\begin{theorem}[Residue Theorem]
Let $f$ be analytic in the region $G$ except for the isolated 
singularities $a_1,a_2,\dots,a_m$. If $\gamma$ is a closed 
rectifiable curve in $G$ which does not pass through any of the 
points $a_k$ and if $\gamma\approx 0$ in $G$, then
\[
  \frac{1}{2\pi i}\int_\gamma\! f = \sum_{k=1}^m 
  n(\gamma;a_k)\Res(f;a_k)\,.
\]
\end{theorem}
\end{document}
